\documentclass[9pt,a4paper]{extarticle}

% ======================
% Codificação, idioma e layout
% ======================
\usepackage[T1]{fontenc}
\usepackage[utf8]{inputenc}
\usepackage[brazil]{babel}
\usepackage{geometry}
\geometry{margin=2cm}
\usepackage{parskip}
\usepackage{setspace}
\usepackage{microtype}
\usepackage{enumitem}

% ======================
% Matemática e símbolos
% ======================
\usepackage{amsmath, amssymb, amsthm, bm}

% ======================
% Figuras e tabelas
% ======================
\usepackage{graphicx}
\usepackage{booktabs}
\usepackage{array}
\usepackage{tabularx}
\usepackage{float}
\usepackage{caption}
\usepackage{subcaption}

% ======================
% Hiperlinks
% ======================
\usepackage{hyperref}
\hypersetup{
  colorlinks=true,
  linkcolor=blue,
  urlcolor=blue
}

% ======================
% Início do documento
% ======================
\begin{document}

\begin{center}
  \Large \textbf{Lista de Exercícios 1 - Introdução à Análise de Algoritmos e Caracterização de Tempos de Execução} \\[4pt]
  \normalsize Disciplina: PCC104 - Projeto de Análise de Algoritmos
\end{center}

\textit{A lista de exercícios a seguir está baseada no livro-texto adotado na disciplina, disponível na biblioteca digital da UFOP. Para acessá-lo, entre no sistema MinhaUFOP e navegue até: Biblioteca Digital → E-books Minha Biblioteca. Em seguida, localize a obra utilizando a referência indicada abaixo:}

\textit{CORMEN, Thomas H.; LEISERSON, Charles E.; RIVEST, Ronald L.; STEIN, Clifford. Algoritmos: teoria e prática. 4. ed. Rio de Janeiro: Grupo GEN | LTC, 2022.}

\section*{Capítulo 1}

\begin{enumerate}[resume]
    \item Exercício 1.2-2
    \item Exercício 1.2-3
    \item Problema 1-1
\end{enumerate}

\section*{Capítulo 2}

\begin{enumerate}[resume]
    \item Exercício 2.1-1
    \item Exercício 2.1-2
    \item Exercício 2.1-3
    \item Exercício 2.1-4
    \item Exercício 2.2-1
    \item Problema 2-2
\end{enumerate}

\section*{Capítulo 3}

\begin{enumerate}[resume]
    \item Exercício 3.2-1
    \item Exercício 3.2-2
    \item Exercício 3.2-3
    \item Exercício 3.2-6
    \item Exercício 3.2-7
    \item Problema 3-1
    \item Problema 3-2
    \item Problema 3-3
    \item Problema 3-4
    \item Problema 3-5
\end{enumerate}

\end{document}

